\section{Intrinsic resolution at arbitrary exponent collisions}\label{sec:collision-geometry}
We now prove the scalar resolution theorem. Maximal Vandermonde products
on the future exponent set determine both attained checkpoint resolution
and causal memory. They are defined at an exact collision without choosing
a path into it. The finite spectral interpretation is
Proposition~\ref{prop:ordinary-spectrum}; the argument here establishes
attainability, the global image bound, and compatibility with causal updates.

\subsection{Calibration chamber and determinant volumes}
Fix $r\ge2$ and a compact subset $K$ of the open chamber
\[
 \{(a_1,\ldots,a_{r-1}):0<a_1<\cdots<a_{r-1}\},
 \qquad a_0=0.
\]
The parameter of the experiment is $t\in[0,1]$, with one fixed
full-support prior $\mu$. The known calibration is $a\in K$.
One fixed matrix $C=(c_{ji})$ of column rank $r$ defines the cells
\[
 k_j^a(t)=\sum_{i=0}^{r-1}c_{ji}t^{a_i},\qquad
 \sum_jk_j^a(t)=1,\qquad k_j^a(t)\ge k_*>0.
\]
The command cube is $[\eta,1-\eta]^J$, $0<\eta<1/2$, as before.
These assumptions are imposed on all of $K$; no regularity or
semialgebraicity of $K$ is required. Compactness inside the chamber
separates the one-step exponents, but \emph{does not} separate their
additive sums. The finitely many report likelihoods have a common
positive lower bound $\kappa$.

Fix $N\ge2$, and choose once and for all
$H\ge\max\{1,N\max_{a\in K}a_{r-1}\}$. For $m\ge1$ write
\[
 \Lambda_m=\{\alpha\in\N_0^r:|\alpha|=m\},\qquad
 \Lambda_m^+=\Lambda_m\setminus\{me_0\},\qquad
 q_m=|\Lambda_m^+|,
 \quad x_\alpha(a)=H^{-1}\sum_i\alpha_i a_i.
\]
Here $\N_0$ includes zero. Formal labels are retained even when their
values coincide. All $x_\alpha$, $\alpha\in\Lambda_m^+$, lie in
$[a_*/H,1]$, where $a_* =\min_{a\in K}a_1>0$.
For $1\le\ell\le q_m$ define
\begin{equation}\label{eq:determinant-volume}
 \mathcal V_{m,\ell}(a)=
 \max_{\substack{J\subset\Lambda_m^+\\|J|=\ell}}
    \prod_{\substack{\alpha,\beta\in J\\\alpha<\beta}}
          |x_\alpha(a)-x_\beta(a)|,
 \qquad \mathcal V_{m,1}=1.
\end{equation}
The ordering inside the product is immaterial. The maximum is zero if
fewer than $\ell$ distinct positive future exponents remain. Put
\begin{equation}\label{eq:intrinsic-profile}
 \begin{split}
 p_{n,m}&=\min\{n(r-1),q_m\},\\
 \Psi_{n,m}(M,a)&=
   \max_{1\le\ell\le p_{n,m}}
        \left(\frac{\mathcal V_{m,\ell}(a)}{M}\right)^{2/\ell},\\
 \Xi_N(M,a)&=\max_{1\le n<N}\Psi_{n,N-n}(M,a).
 \end{split}
\end{equation}
Set $\Psi_{0,m}=\Psi_{n,0}=0$. Changing $H$ to any other fixed
admissible value changes only comparison constants. In particular the
invariant does not depend on a convention for measuring exponents.

Use the failure-factor basis $1/2$ and $1/2+\delta t^{a_i}$,
$1\le i<r$, realized by the fixed right inverse of $C$ for one
sufficiently small $\delta>0$. Its ordered $m$-fold products are the
query menu, sampled uniformly after the memory index is formed. Expansion
in the symbolic variables $X_0,\ldots,X_{r-1}$ gives a fixed matrix
$G_m$ of column rank $|\Lambda_m|$: an invertible linear change of
variables on homogeneous polynomials of degree $m$ is invertible.
Consequently the physical query norm and the norm of the formal raw
moment vector are uniformly equivalent. Coincident exponent values impose
relations on attainable vectors, not a loss of column rank of this
\emph{formal} coefficient matrix.

\subsection{A uniformly conditioned flag}
Order the positive nodes greedily. Choose the first label by a fixed rule;
for example take its least node, breaking ties by a fixed label order.
Having selected $x_1,\ldots,x_{j-1}$, select a remaining label maximizing
\[
 \prod_{i<j}|x-x_i|,
 \qquad d_1=1,\qquad d_j=\prod_{i<j}|x_j-x_i|.
\]
Break every tie by the same label rule. After all distinct values have
been selected, continue through the remaining labels; their $d_j$ are
zero. This is finite Leja ordering. The ordering and its use in Newton
interpolation are classical \cite{Leja,BosEtAl}; the estimates needed
here are proved below, including their zero-pivot form.

\begin{lemma}[Newton scales and maximal determinant volumes]\label{lem:leja-scales}
For every calibration,
$1=d_1\ge d_2\ge\cdots\ge d_{q_m}\ge0$. There is a matrix
$L_a$, whose entries have absolute value at most one, such that
\begin{equation}\label{eq:leja-physical}
 (t^{Hx_i})_{i=1}^{q_m}
     =L_a\operatorname{diag}(d_j)(g_j^a(t))_{j=1}^{q_m},
 \qquad
 g_j^a(t)=[x_1,\ldots,x_j](z\mapsto t^{Hz}).
\end{equation}
Divided differences at repeated nodes have their Hermite values. If
$s$ is the number of distinct positive nodes, the first $s$ rows and
columns of $L_a$ form a lower triangular matrix with diagonal entries
$\pm1$; its inverse has a bound depending only on $q_m$. Thus, for
any vector $v$,
\begin{equation}\label{eq:leja-metric}
 c_m\|\operatorname{diag}(d_j)v\|
 \le\|L_a\operatorname{diag}(d_j)v\|
 \le C_m\|\operatorname{diag}(d_j)v\|.
\end{equation}
Moreover, for every $1\le\ell\le q_m$,
\begin{equation}\label{eq:leja-volume}
 \prod_{j=1}^{\ell}d_j\ \le\ \mathcal V_{m,\ell}(a)
       \ \le\ \ell!\prod_{j=1}^{\ell}d_j.
\end{equation}
All these assertions hold at exact coincidences.
\end{lemma}
\begin{proof}
Every distance is at most one. The maximization defining $d_j$ and the
removal of a candidate imply $d_{j+1}\le d_j$. Put
$P_{j-1}(z)=\prod_{i<j}(z-x_i)$. For $d_j>0$ set
$L_{ij}=P_{j-1}(x_i)/d_j$; for $d_j=0$ set that column to zero.
A previously selected node gives zero when $i<j$, and the greedy choice gives $|L_{ij}|\le1$ on all remaining
nodes. The leading nonzero block is lower triangular with diagonal
$\pm1$. Back substitution bounds its inverse in a fixed dimension.
If $d_j=0$, all remaining node values already occur among the earlier
ones, so $P_{j-1}$ vanishes on every node. Newton's interpolation
identity, evaluated at the nodes, is therefore exactly
\eqref{eq:leja-physical}, also when there are repeats. Only the first
$s$ scaled coordinates can be nonzero, proving \eqref{eq:leja-metric}.

For $\ell\le s$, evaluation of the monic polynomials
$P_0,\ldots,P_{\ell-1}$ on any $\ell$ chosen nodes has determinant
equal to their Vandermonde determinant. Its matrix is the selected rows
of $L_{:,1:\ell}$ times $\operatorname{diag}(d_1,\ldots,d_\ell)$.
The Leibniz formula bounds the first determinant by $\ell!$.
For the first $\ell$ selected nodes its absolute value is one, giving
both bounds in \eqref{eq:leja-volume}. For $\ell>s$, both sides
vanish. No division by a zero pivot has been made.
\end{proof}

The next assertion is the attainable part of the argument. Bounds for a
Vandermonde matrix alone would not prove it.

\begin{lemma}[Uniform attainable Newton flags]\label{lem:newton-attainment}
For every $n,m\ge1$ with $n+m\le N$, let
$z_a(h)=(\E[g_j^a(t)\mid h])_{j=1}^{q_m}$ and
$p=p_{n,m}$. At one common interior all-failure command tuple, the first
$p$ coordinates of $z_a$ have a surjective differential. Its least row
singular value is bounded below uniformly over $a\in K$ and every
ordering of the formal nodes, including limiting orderings at collisions.
Under independent uniform acquisition commands, their all-failure
subprobability law dominates $\beta$ times uniform probability on a
translated $p$-cube of side $2\rho$, where $\beta,\rho>0$ are uniform.
\end{lemma}
\begin{proof}
We first justify continuity at all repeated-node configurations, without
an assumption on how a collision is approached. The Hermite--Genocchi
formula \cite{deBoor} expresses $g_j^a(t)$ as the integral of
$(H\log t)^{j-1}t^{Hz}$ over the simplex of convex combinations of
$x_1,\ldots,x_j$; the simplex has volume $1/(j-1)!$.
Since $Hz\ge a_*>0$,
\[
 |g_j^a(t)|\le
  \frac{H^{j-1}}{(j-1)!}\,t^{a_*}|\log t|^{j-1},
 \qquad 0<t\le1.
\]
The right side is bounded and tends to zero at $t=0$. The same argument
with higher logarithmic powers gives continuity, uniform in $t$, as the
nodes vary in their compact range. Define $g_j^a(0)=0$.

For any initial multiset of $p$ nodes, the $p$ divided-difference
functionals are independent on polynomials of degree at most $p-1$:
their values on $1,z,\ldots,z^{p-1}$ form an upper triangular matrix
with diagonal one. They span the Hermite evaluation functionals at the
distinct nodes, with all derivative orders below the corresponding
multiplicities. Applied to $t^{Hz}$, their span is therefore the complete
exponential-polynomial space
\[
 \spanop\{t^{H\lambda}(\log t)^k:
          \lambda\text{ in the multiset},\ 0\le k<m_\lambda\}.
\]
This argument permits repeated nodes at nonadjacent positions. Adjoin the
constant function. The resulting $p+1$ tests form a complete confluent
Chebyshev space, not a selection of isolated high derivatives.

Choose the common binomial factors
$F_i=(1+c_it^{a_{r-1}})/2$, with distinct sufficiently small positive
$c_i$. Their command coefficients are independent of $a$. By
Lemma~\ref{lem:binomial-tangent} their product tangent has
$n(r-1)+1$ distinct monomials. The mixed and confluent pairing of
Lemmas~\ref{lem:mixed-moment} and \ref{lem:confluent-positive} has
rank $p+1$ against the selected tests, at every node configuration.
In the normalized differential
\[
 Q\longmapsto(\mu F)^{-1}(LQ-v\,e_0LQ),
 \qquad v=LF/\mu F,
\]
the product belongs to the tangent and $e_0v=1$. Thus exactly the
constant direction is lost and the differential onto the other $p$
coordinates is surjective.

For each fixed permutation of formal labels this derivative is continuous
on $K$, including coincidences, by the uniform bounds just proved and
positivity of evidence. There are finitely many permutations. Compactness
therefore supplies a common positive least row singular value. First and
second derivatives in command coordinates are uniformly bounded because
the numerator is a polynomial and the evidence is at least $\kappa^n$.

To record the probability step explicitly, complete the derivative's row
space by orthonormal kernel coordinates. A ball of uniform radius in
command space gives a square map whose derivative differs from its
invertible central derivative by at most half its least singular value.
The contraction proof of the inverse theorem gives a fixed-size product
box in output and kernel coordinates, with its inverse inside the command
cube. The inverse Jacobian determinant is bounded below by the reciprocal
of a uniform power of the forward derivative bound. The command density
is constant, and the joint density of those commands and the all-failure
word is at least that density times $\eta^n$. Integrating over the
kernel box proves the stated subprobability minorization. No prior
density is involved and no failure evidence is discarded.
\end{proof}

\subsection{The global attainable law}
\begin{theorem}[Collision-uniform checkpoint classification]\label{thm:intrinsic-checkpoint}
Under the assumptions of this section there are $0<c\le C<\infty$
such that, for every $a\in K$ and integer $M\ge1$,
\begin{equation}\label{eq:intrinsic-checkpoint}
 c\Psi_{n,m}(M,a)\le\overline R_{M;n,m}^{a}
 \le R_{M;n,m}^{a,\max}\le C\Psi_{n,m}(M,a).
\end{equation}
The unconditional lower law uses independent uniform acquisition commands.
The upper code is deterministic, has at most $M$ reachable representatives,
and works on every admitted history. The lower law permits the independent
coding randomization of Definition~\ref{def:finite-state}. Constants may
depend on $K,C,\eta,\mu,n,m$, but not on $a$ or $M$.
\end{theorem}
\begin{proof}
Let $D_a=\operatorname{diag}(d_j)$. For one report word each coordinate
of $z_a(h)$ has the form $N_j(u)/Z(u)$, with numerator and evidence
polynomials of degree at most $n$ in the command variables and
$Z(u)\ge\kappa^n$. Thus
\[
 S_a=\{D_az_a(h):|h|=n\}
\]
is compact and has coefficient-independent bounded semialgebraic format:
it is a finite union of projections of the formulas
$Z(u)y_j=d_jN_j(u)$ on the command cube. Prior integrals and node
values are real coefficients in these formulas. They need not depend
semialgebraically on $a$.

Individually normalizing the $n$ report factors puts them in affine
hyperplanes of dimension $r-1$. The moment map is rational in those
coordinates. Hence $\dim S_a\le p$. Lemma~\ref{lem:newton-attainment}
also gives $|z_{a,j}|\le B$, so $S_a$ lies in the rectangle with
half-widths $Bd_1\ge\cdots\ge Bd_{q_m}$. By
Lemma~\ref{lem:tame-rectangle}, it has at most $M$ reachable
representatives with covering radius bounded by
\[
 C\max_{\ell\le p}
       \left(M^{-1}\prod_{j\le\ell}d_j\right)^{1/\ell}.
\]
The physical metric equivalence follows from \eqref{eq:leja-metric}
and the full-column-rank product-probe matrix $G_m$. A fixed tie rule
makes nearest-representative encoding Borel. Squaring and applying
\eqref{eq:leja-volume} gives the upper bound in
\eqref{eq:intrinsic-checkpoint}, with the integer budget already enforced
by Lemma~\ref{lem:tame-rectangle}.

For the lower bound, first average any randomized decoder conditional on
its label; squared loss cannot increase. Orthogonally project its query
centers onto the affine physical query span, apply a fixed left inverse
of $G_m$, and then the inverse of the leading nonzero block of $L_a$.
These maps have uniformly bounded norms. Projection onto any first
$\ell\le\min(p,s)$ coordinates reduces the squared-error problem to
the marginal of the cube in Lemma~\ref{lem:newton-attainment}, scaled
by $d_1,\ldots,d_\ell$. The probability that $M$ radius-$b$ balls
cover this uniform rectangle is at most
\[
 C_\ell M b^\ell/(d_1\cdots d_\ell).
\]
Choose $b=c_\ell(d_1\cdots d_\ell/M)^{1/\ell}$ with the
constant small enough that at least half the minorized mass remains.
Its mean squared error is at least a fixed multiple of $b^2$.
An encoder's randomized assignment cannot improve upon the distance to
the nearest of its centers. Independent public coding randomness can be
fixed and then averaged, since it does not generate the acquisition law.
Maximize over $\ell$. For $\ell>s$ the determinant volume is zero,
so no additional term is needed. Equation~\eqref{eq:leja-volume}
now gives the lower bound. The unconditional optimum is bounded above
by the minimax optimum, completing the proof.
\end{proof}

\begin{theorem}[Causal classification across all collision strata]\label{thm:intrinsic-streaming}
For the same fixed experiment and horizon, there exist $c,C>0$ such that
for every known $a\in K$ and every integer $M\ge1$ there exists a
\emph{single, stage-compatible} deterministic $M$-state transducer with
maximum-history, maximum-checkpoint regret at most $C\Xi_N(M,a)$.
Every $M$-state transducer has maximum unconditional checkpoint regret at
least $c\Xi_N(M,a)$ under one common independent uniform exploration law.
Equivalently,
\begin{equation}\label{eq:intrinsic-streaming}
 c\Xi_N(M,a)\le\mathfrak R_{M,N}^{a,\mathrm{av}}
 \le\mathfrak R_{M,N}^{a,\max}\le C\Xi_N(M,a).
\end{equation}
The transducer may depend on $a$ and $M$. Constants are uniform in both;
codebooks are not asserted to be calibration-blind. It stores only its
representative index. The clock, known calibration and fixed real-valued
program are read-only, as in Definition~\ref{def:finite-state}.
\end{theorem}
\begin{proof}
At stage $n$ let $m=N-n$ and store a representative in the reachable
\emph{raw} moment image indexed by $\Lambda_m$, including the constant
coordinate $v_{me_0}=1$. For a new factor
$f(t)=\sum_i f_i(g,x)t^{a_i}$, the exact update is
\begin{equation}\label{eq:formal-multiindex-update}
 v'_\alpha=
 \frac{\sum_i f_i(g,x)v_{\alpha+e_i}}
      {\sum_i f_i(g,x)v_{e_i+(m-1)e_0}},
 \qquad \alpha\in\Lambda_{m-1}.
\end{equation}
Every index on the right belongs to $\Lambda_m$. Coincident labels
have equal moment values and cause no ambiguity. On the segment joining
any two reachable moment vectors, the denominator is a report probability
under a posterior mixture and is at least $\kappa$. Coefficients are
uniformly bounded and raw moments belong to $[0,1]$. Differentiating the
quotient gives a Lipschitz constant independent of all exponent gaps.

At every stage Theorem~\ref{thm:intrinsic-checkpoint} provides a raw
reachable codebook of cardinality at most $M$ and covering radius at most
$C_n\Psi_{n,N-n}(M,a)^{1/2}$. Update a representative by
\eqref{eq:formal-multiindex-update}, then quantize in the next reachable
codebook. Updating a reachable representative produces a reachable next
state, so this quantizer is always evaluated on its domain. If $e_n$
is the raw error, then
\[
 e_0=0,\qquad e_{n+1}\le L_ne_n+
        C_{n+1}\Psi_{n+1,N-n-1}(M,a)^{1/2}.
\]
The finite recurrence bounds every $e_n$ by
$C_N\Xi_N(M,a)^{1/2}$. The query matrix converts this to the
asserted squared-regret bound. The exact prefix appears only in this
error analysis, not in the transducer. No divided-difference coordinate,
reciprocal pivot or stored command tape is used in a transition.

Conversely, at each checkpoint the streaming index is an $M$-message
encoder of the prefix. All command laws used in
Theorem~\ref{thm:intrinsic-checkpoint} are prefixes of one common
exploration law. Each transducer therefore satisfies every checkpoint
lower bound; taking their maximum and then the infimum proves the lower
bound in \eqref{eq:intrinsic-streaming}. Finitely many checkpoint
constants can be replaced by their positive minimum. One checkpoint
and its remaining query trials are executed in a run, as before.
\end{proof}

\begin{corollary}[Intrinsic bit law]\label{cor:intrinsic-bits}
Let $B_*(\varepsilon,a)$ be the least integer persistent bit budget
for maximum-history streaming regret at most $\varepsilon$ in this
experiment. Uniformly for $0<\varepsilon\le1$ and $a\in K$,
\begin{equation}\label{eq:intrinsic-bits}
 B_*(\varepsilon,a)=
 \max_{\substack{1\le n<N\\1\le\ell\le p_{n,N-n}}}
 \left[\frac{\ell}{2}\log_2\frac1\varepsilon+
       \log_2\mathcal V_{N-n,\ell}(a)\right]_++O(1).
\end{equation}
A zero determinant volume contributes zero to this maximum, by interpreting
its logarithm as $-\infty$ before taking the positive part. The same law
holds for the unconditional maximum-checkpoint criterion. At each fixed
calibration the small-error slope is $D_A(N)/2$.
\end{corollary}
\begin{proof}
With $M=2^B$, solving each inequality
$(\mathcal V_{m,\ell}/M)^{2/\ell}\le\varepsilon$ gives the displayed
lower bound on $B$. Constants in \eqref{eq:intrinsic-streaming} change
it by at most a bounded amount, since $\ell$ is bounded. Integer rounding
costs at most one further bit. For fixed $a$, the largest nonzero index
at checkpoint $(n,m)$ is
$\min\{n(r-1),|mA|-1\}=d_A(n,m)$, proving the slope assertion.
\end{proof}

\subsection{Collision orders and a finite tree formula}
The determinant form has a consequence not requiring interpolation
coordinates: it classifies arbitrary orders of tangency to a collision.
It also gives a finite optimization rule for all resolution phases.

\begin{theorem}[Collision-tree determination of memory phases]\label{thm:collision-tree}
Let $a(\theta)\in K$, $0\le\theta\le\theta_0$, be a known path.
For a fixed $m$, identify formal sums that are identical along the path.
Suppose that for every remaining pair of labels
\[
 |x_\alpha(\theta)-x_\beta(\theta)|
       \asymp\theta^{w_{\alpha\beta}},\qquad
 w_{\alpha\beta}\in[0,\infty),\quad 0<\theta\le\theta_0.
\]
After decreasing $\theta_0$ this hypothesis holds for every real analytic
path whose nonidentical sums are analytic at zero. Define
\begin{equation}\label{eq:tree-energy}
 \beta_{m,\ell}=
   \min_{\substack{J\text{ a set of remaining labels}\\|J|=\ell}}
         \sum_{\{\alpha,\beta\}\subset J}w_{\alpha\beta},
 \qquad \beta_{m,1}=0.
\end{equation}
Then $\mathcal V_{m,\ell}(a(\theta))\asymp
\theta^{\beta_{m,\ell}}$. For unavailable cardinalities the volume is
zero. Consequently the checkpoint law is
\begin{equation}\label{eq:tree-law}
 \overline R_{M;n,m}^{a(\theta)}\asymp
 \max_{\ell\le p_{n,m}}
       \theta^{2\beta_{m,\ell}/\ell}M^{-2/\ell},
\end{equation}
with zero-volume terms omitted; the streaming law is its maximum over
checkpoints. The exponents $\beta_{m,\ell}$ depend only on the pairwise
collision orders, and can be computed by a finite tree allocation rule.
\end{theorem}
\begin{proof}
Each fixed subset product has order equal to the sum of its pair orders.
There are finitely many subsets, so their maximum has the smallest such
order. This proves the determinant assertion and then
\eqref{eq:tree-law} by Theorems~\ref{thm:intrinsic-checkpoint} and
\ref{thm:intrinsic-streaming}. Constants may depend on the path's
nonzero leading coefficients; the intrinsic determinant theorem, in
contrast, is already uniform on $K$.

For completeness we give the allocation rule and its justification. The
triangle inequality for node differences implies
\[
 w_{\alpha\gamma}\ge
       \min\{w_{\alpha\beta},w_{\beta\gamma}\}.
\]
Indeed a smaller order on the left would contradict the upper bound by
the sum of the two differences on the right as $\theta\downarrow0$.
Thus thresholds on $w$ define nested equivalence classes. Make a rooted
tree from these classes, with root height zero; an internal cluster $C$
has height $h_C=\min_{\alpha\ne\beta\in C}w_{\alpha\beta}$,
and its children are the equivalence classes for orders strictly greater
than $h_C$. Split the root at orders greater than zero. Singleton
children are leaves. Redundant equal-height vertices can be omitted.
For an internal nonroot vertex put
$\Delta_C=h_C-h_{\operatorname{parent}(C)}$. The sum of these increments
on the common ancestor chain of two leaves is exactly their pair order.
Hence, for a selected leaf set $J$,
\begin{equation}\label{eq:tree-decomposition}
 \sum_{\{\alpha,\beta\}\subset J}w_{\alpha\beta}
   =\sum_{C\text{ internal},\ C\ne\mathrm{root}}
               \Delta_C\binom{|J\cap C|}{2}.
\end{equation}
Set $F_C(0)=F_C(1)=0$ for a leaf and $F_C(k)=+\infty$ for $k>1$.
For an internal vertex, including the root with $\Delta_{\rm root}=0$,
compute
\begin{equation}\label{eq:tree-dp}
 F_C(k)=\Delta_C\binom{k}{2}+
   \min_{\substack{\sum_{D\text{ child of }C}k_D=k}}
                 \sum_{D\text{ child of }C}F_D(k_D).
\end{equation}
Equation~\eqref{eq:tree-decomposition} proves this recurrence by induction.
It gives $F_{\rm root}(\ell)=\beta_{m,\ell}$. Successive finite
min-convolutions compute the table; no optimization over measures,
posterior encoders or codebooks is involved. For analytic paths the pair
orders are the finite orders of their first nonzero Taylor coefficients,
after identical sums have been identified. This proves the assertion.
\end{proof}

For the affine family of Section~\ref{sec:confluence}, different limiting
clusters have pair order zero, and distinct members of a cluster have pair
order one. Thus a choice of $k$ nodes from that cluster costs
$\binom{k}{2}$; its successive costs are $0,1,\ldots,s-1$.
Minimizing their sum picks the smallest costs from all clusters. Therefore
$\beta_{m,\ell}=\nu_{m,1}+\cdots+\nu_{m,\ell}$. This recovers
Theorem~\ref{thm:attainable-checkpoint}, while its separate proof remains
available. Higher orders of contact between distinct branches require the
pair orders above and are not encoded by the affine multiplicity list.

\subsection{Intersecting collision loci and arbitrary tangency order}
Here is a two-parameter consequence with different ranks on intersecting
loci and arbitrarily deep resolution phases along tangent approaches.
It illustrates the geometry of the determinant theorem rather than
requiring a separate asymptotic analysis for each path.

\begin{corollary}[A two-parameter collision arrangement]\label{cor:two-parameter}
Let $A_{u,v}=\{0,1,2+u,3+v\}$, with
$|u|,|v|\le1/16$, and take the four cells
\[
 k_0(t)=\tfrac14+\tfrac1{16}(t+t^{2+u}+t^{3+v}),\qquad
 k_i(t)=\tfrac14-\tfrac1{16}t^{a_i},\quad 1\le i\le3.
\]
Use any fixed admitted command cube and a fixed full-support prior.
For total horizon five set
$\rho=\max\{|u|,|v|\}$ and
$\tau=\min\{|u|,|v-u|,|v-2u|\}$. Then, uniformly on the square,
\begin{equation}\label{eq:two-parameter-law}
 \mathfrak R_{M,5}^{u,v}\asymp
 \max\left\{M^{-1/3},\ \rho^{1/2}M^{-1/4},\
           (\rho^2\tau)^{2/9}M^{-2/9}\right\}
\end{equation}
for either streaming criterion. The exact peak dimension is nine off the
three collision lines, eight on any such line away from their intersection,
and six at $(0,0)$.

Along $u=\theta$, $v=\theta+\theta^k$ with an integer $k\ge2$,
the three regimes have respective orders
\[
 M^{-1/3},\qquad \theta^{1/2}M^{-1/4},\qquad
       \theta^{(4+2k)/9}M^{-2/9}.
\]
Their crossover budgets have orders $\theta^{-6}$ and
$\theta^{-(8k-2)}$, with regrets of orders $\theta^2$ and
$\theta^{2k}$, respectively. These are comparison orders, not exact
integer thresholds.
\end{corollary}
\begin{proof}
The cells sum to one, lie between $3/16$ and $7/16$, and their coefficient
matrix has rank four: its last three rows separate the three nonconstant
coefficients and the sum of its rows gives the constant. They therefore
satisfy the compact-chamber hypotheses.

At checkpoint $n=3,m=2$ the nine positive formal exponents are
\[
 1,\quad 2,\ 2+u,\quad 3+u,\ 3+v,\quad
 4+2u,\ 4+v,\quad 5+u+v,\quad 6+2v.
\]
They form six uniformly separated groups, with three within-group gaps
$|u|,|v-u|,|v-2u|$. Denote these gaps in decreasing order by
$g_1\ge g_2\ge g_3$. For $\ell\le6$ the determinant volumes
are comparable to one. For $\ell=6+j$, $1\le j\le3$, they are
comparable to $g_1\cdots g_j$: choose one exponent from every group
and a second from the $j$ groups with largest gaps for the lower bound.
Any selection of $6+j$ exponents must double at least $j$ groups;
all other factors are bounded, proving the upper bound. Fixed powers of
$H$ only change constants. Since $p_{3,2}=9$, its profile is
\[
 \max\{M^{-1/3},\ g_1^{2/7}M^{-2/7},\
              (g_1g_2)^{1/4}M^{-1/4},\
              (g_1g_2g_3)^{2/9}M^{-2/9}\}.
\]
The signed differences satisfy $(v-u)-(v-2u)=u$, so the largest
absolute gap is at most twice the second largest. Moreover
$g_1\asymp\rho$, hence $g_2\asymp\rho$ and $g_3=\tau$.
The seven-dimensional term is the geometric interpolation, with weights
$3/7$ and $4/7$, of the six- and eight-dimensional terms, up to uniform
constants. This proves the expression in \eqref{eq:two-parameter-law}
for the checkpoint, also on the collision lines.

At the other checkpoints, the attainable dimensions are bounded by three,
six and three, respectively. Since every normalized gap is at most one,
their squared covering errors are at most $CM^{-1/3}$. The causal
classification therefore gives the same entire-horizon profile and the
checkpoint-three lower bound supplies its converse. Counting distinct
two-fold exponents gives the three exact peak dimensions asserted.

Finally, on the stated path $\rho\asymp\theta$ and
$\tau\asymp\theta^k$. Substitution gives the three displayed terms.
Equality in order of the first two gives $M=\theta^{-6}$, at regret
$\theta^2$; equality in order of the last two gives
$M=\theta^{-(8k-2)}$, at regret $\theta^{2k}$. As $k>1$, these
budgets are strictly separated as $\theta\downarrow0$ and all three
regimes occur. The exponent $k$ records contact with a collision locus,
not an additional assumed observable jet.
\end{proof}
